% =====================================================================
\section{领头区域：初态与末态相互作用}\label{sec:dy-regions}
% =====================================================================

DY 的约化图（\cref{fig:reduced-dy}）与 DIS、$e^+e^-$ 的约化图形式相同：硬子图、两个共线子图与一个软子图。差别在于软胶子所连接的线的“时间性质”。在 $e^+e^-$ 湮灭中，全部共线线都属于末态；在 DIS 中，靶子图 $J_A$ 既含初态线（硬顶点之前的活动部分子）也含末态线（飞向末态的旁观者），但另一个共线子图——被击出的夸克喷注——只含末态线，于是一个软动量的 $k^+$ 分量只看到同一半平面的极点，围道可以变形。在 DY 中，$J_A$ 与 $J_B$ 都同时含有初态线与末态线：两个强子的旁观者之间交换的软胶子在 $k^-$ 上被 $J_A$ 钳住，在 $k^+$ 上被 $J_B$ 钳住。这就是 DY 特有的 Glauber 问题；本节先把领头区域列出来，再在单圈上把钳夹显式算出。

\begin{proposition}[DY 的领头区域 \tagA]\label{prop:dy-regions}
在 Feynman 规范中，单举 DY 截面（对全部割图求和）的领头幂次区域由\cref{thm:ir-leading-regions}（全阶论证，\tagA）的约化图给出（\cref{fig:reduced-dy}）：
\begin{enumerate}[label=(\roman*),leftmargin=2.2em,itemsep=1pt,topsep=2pt]
\item 末态割线两侧各一个硬子图 $H$、$H^\dagger$，其中的线离壳 $\order{Q^2}$；
\item 共线子图 $J_A$、$J_B$，各自含强子顶点与全部旁观者并穿过末态割线，在割线每一侧经一个物理部分子（夸克的好分量或横向胶子）加任意多个纵向极化的共线胶子连到 $H$（$H^\dagger$）；
\item 软子图 $S$，只经胶子连到 $J_A$、$J_B$，不连到 $H$、$H^\dagger$；
\item 作为 (iii) 的子区域，Glauber 动量 $k\sim(\lambda^2,\lambda^2,\lambda)Q$（\cref{def:ir-glauber}），特别是两个强子的旁观者之间的交换，它们与软动量一样是领头幂次。
\end{enumerate}
在单个割图中，$J_A$、$J_B$ 都既含硬顶点之前的初态线，也含硬顶点之后、割线之前的末态线；旁观者的软相互作用因而既有初态相互作用，也有末态相互作用。（本命题只涉及部分子态。）
% errtype: E-S
\uses{thm:ir-leading-regions,def:dy-kin,def:ir-glauber}
\end{proposition}

\begin{proofsketch}
(1) 由 Landau 方程与 Coleman--Norton 图像得到钳夹面，由幂次计数挑出 $\omega=0$ 的那些，得到\cref{thm:ir-leading-regions}的约化图；对 DY 该定理已把旁观者之间的 Glauber 子区域列入（\cite[第~5.8--5.11~节]{Collins2011}，\cite[第~8.7~节]{CollinsSoperSterman1989}）。
(2) Glauber 区确为领头幂次：相对于不含该胶子的图，测度 $\dd k^+\dd k^-\dd^2\mathbf k_T\sim\lambda^6Q^4$，胶子传播子 $\simeq-1/\mathbf k_T^2\sim\lambda^{-2}Q^{-2}$，两条新的旁观者传播子各离壳 $\order{\lambda^2Q^2}$，共 $\lambda^{-4}Q^{-4}$，分子中与大分量的耦合给出 $(2p_A^+)(2p_B^-)\sim Q^2$，乘积为 $\lambda^0$，与 $Q$ 无关\cite[第~5.6.2~节]{Collins2011}（此步只给出 E-S 型陈述）。
(3) 旁观者穿过割线，软胶子在割线同一侧既可附着在硬顶点 $V_H$ 之前的旁观者段，也可附着在其后的段上，于是两种相互作用在同一个割图中出现。
完整论证见\cite[第~5~章与第~14.3~节]{Collins2011}。
\end{proofsketch}

\paragraph{地位与缺口}\tagA。最弱的一步是 (1)–(2) 中的幂次计数：它们是标度陈述（E-S），积分范数下的形式是\cref{asm:appN-envelope}、\cref{asm:appN-separation}（\tagA），在所有阶都没有以积分形式写出。Feynman 规范中领头区域还不完备：额外的不连通硬散射只在对图求和后相消\cite{Labastida:1984gy}，Collins 把这一点列为 DY 证明中“尚存的问题”之一\cite[第~14.4.8~节]{Collins2011}。对强子态，“微扰论找到的领头区域就是全部领头区域”是桥接假设\cref{asm:intro-bridge}(Br1)，不属于本命题。

\cref{fig:glauber-exchange}把 DY 中旁观者之间的初态 Glauber 交换与 DIS 中被击出夸克与靶残余之间的末态交换并列。图中说明了单个割图中极点分居两侧、围道被钳住的情形；需要补充的是：离开 Glauber 区须沿 $e_G$ 同时形变两个分量（\cref{fig:method-glauber-direction}）——只推动 $k^-$ 一个分量，即使推到 $|k^-|\sim|\mathbf k_T|$，也仍有 $|k^+k^-|\sim\lambda^3\ll\mathbf k_T^2$，动量并没有离开 Glauber 区（\cref{rem:method-two-component}）。

%%% BEGIN fig:glauber-exchange
\begin{figure}[htbp]
\centering
\begin{tikzpicture}[fd,
  t2 hadron/.style={draw=black!70,line width=2.2pt,postaction={decorate,decoration={markings,
      mark=at position 0.6 with {\arrow[black!70]{Stealth[length=7pt,width=7.5pt]}}}}},
  t2 lab/.style={font=\footnotesize},
  t2 time/.style={lcaxis,gray!60!black}]
  % ======================= (a) Drell-Yan: initial-state exchange between spectators =======================
  \begin{scope}
    \node[jet blob,circle,minimum size=0.66cm,inner sep=0pt] (A) at (-2.45,1.0) {$A$};
    \node[jetB blob,circle,minimum size=0.66cm,inner sep=0pt] (B) at (-2.45,-1.25) {$B$};
    \node[hard blob,circle,minimum size=0.6cm,inner sep=0pt] (H) at (1.65,1.5) {$H$};
    \draw[t2 hadron] (-3.6,1.0) -- (A.west);
    \draw[t2 hadron] (-3.6,-1.25) -- (B.west);
    % active partons: quark from A, antiquark from B (fermion-number arrow points back into B)
    \draw[quark] (A.12) -- (H.185);
    \draw[quark=0.45] (H.235) -- (B.42);
    \draw[photon] (H.east) -- (2.85,1.5);
    % spectators (remnants) of A and B, continuing into the final state
    \coordinate (SA) at (A.-28); \coordinate (XA) at (2.8,0.25);
    \coordinate (SB) at (B.-8);  \coordinate (XB) at (2.8,-1.2);
    \draw[quark=0.25] (SA) -- (XA);
    \draw[quark=0.25] (SB) -- (XB);
    % Glauber exchange before the hard collision (left of H)
    \coordinate (Gt) at (0.78,2.0); \coordinate (Gb) at (0.78,-2.0);
    \coordinate (GA) at (intersection of SA--XA and Gt--Gb);
    \coordinate (GB) at (intersection of SB--XB and Gt--Gb);
    \draw[glauber] (GA) -- (GB);
    \node[vertex] at (GA) {}; \node[vertex] at (GB) {};
    \node[t2 lab,text=red!55!black,anchor=west] at ($(GA)!0.5!(GB)+(0.16,0)$) {Glauber};
    \node[t2 lab,anchor=west] at (2.85,1.5) {$\gamma^*$};
    \node[t2 lab,anchor=west] at (XA) {$X_A$};
    \node[t2 lab,anchor=west] at (XB) {$X_B$};
    \node[t2 lab,anchor=south] at (-3.25,1.1) {$P_A$};
    \node[t2 lab,anchor=north] at (-3.25,-1.35) {$P_B$};
    \draw[t2 time] (-3.5,-1.95) -- (3.15,-1.95) node[above left,t2 lab,text=gray!50!black] {时间};
    \node at (-0.2,-2.5) {(a)\ DY：旁观者之间的初态相互作用};
  \end{scope}
  % ======================= (b) DIS: final-state exchange between struck quark and remnant =======================
  \begin{scope}[xshift=7.6cm]
    \node[jet blob,circle,minimum size=0.66cm,inner sep=0pt] (A) at (-2.0,-0.95) {$A$};
    \node[hard blob,circle,minimum size=0.6cm,inner sep=0pt] (H) at (0.15,0.45) {$H$};
    \draw[t2 hadron] (-3.5,-0.95) -- (A.west);
    \draw[photon] (-1.35,1.75) -- (H.135);
    \draw[quark] (A.35) -- (H.215);
    \draw[quark] (H.20) -- (2.45,1.45);
    \draw[quark=0.42] (A.-15) -- (2.45,-1.25);
    \coordinate (Gq) at ($(H.20)!0.52!(2.45,1.45)$);
    \coordinate (Gr) at ($(A.-15)!0.735!(2.45,-1.25)$);
    \draw[glauber] (Gq) -- (Gr);
    \node[vertex] at (Gq) {}; \node[vertex] at (Gr) {};
    \path[mom={$q$},mom sep=0.28cm] (-1.35,1.75) -- (H.135);
    \node[t2 lab,anchor=west] at (2.45,1.45) {被击出的夸克};
    \node[t2 lab,anchor=west] at (2.45,-1.25) {靶残余};
    \node[t2 lab,text=red!55!black,anchor=west] at ($(Gq)!0.5!(Gr)+(0.17,0)$) {Glauber};
    \node[t2 lab,anchor=south] at (-3.15,-0.85) {$P$};
    \draw[t2 time] (-3.5,-1.95) -- (2.7,-1.95) node[above left,t2 lab,text=gray!50!black] {时间};
    \node at (-0.45,-2.5) {(b)\ DIS：被击出夸克与靶残余之间的末态相互作用};
  \end{scope}
\end{tikzpicture}
\caption{Glauber 胶子交换（暗红色）。时间自左向右，硬碰撞 $H$ 之前的交换为初态相互作用，之后为末态相互作用。
(a)~Drell--Yan：强子 $A$、$B$ 的旁观者（强子残余，进入末态 $X_A$、$X_B$）在硬碰撞之前交换 Glauber 胶子（$|k^+k^-|\ll k_T^2$）；来自 $B$ 的入射反夸克线与 $A$ 的旁观者线的交叉不是顶点。
在单个割图中，$A$ 一侧（旁观者线与活动部分子线）的 $k^-$ 极点与 $B$ 一侧的 $k^+$ 极点都分居实轴两侧，围道被钳在 Glauber 区（\cref{fig:glauber-contour}），只有对旁观者的末态切割求和（幺正性）之后才能把围道变形出 Glauber 区；
(b)~DIS：被击出的夸克与靶残余在硬碰撞之后交换 Glauber 胶子；对单举 DIS 它们由幺正性相消，而在 SIDIS 的 TMD 中它们组合成未来指向的 Wilson 线，
与 Drell--Yan 中由初态相互作用产生的过去指向 Wilson 线相对照（Sivers 函数因而在两过程间变号）。}
\label{fig:glauber-exchange}
\end{figure}
%%% END fig:glauber-exchange

钳夹的机制可以在单圈上完全显式地看到。下面的命题把\cref{prop:ir-glauber-trapping}（\tagR，$R_{\mathrm{PT}}$）的一般判据用到 DY 的旁观者上；同一个单圈例子见\cite[第~5.6.2~节]{Collins2011}与\cite[第~9.3~节，式~(166)]{CollinsSoperSterman1989}。

\begin{proposition}[单个割图中的 Glauber 钳夹 \tagR]\label{prop:dy-glauber-pinch}
（$R_{\mathrm{PT}}$：在重整化微扰论中对红外正规化的部分子态严格。）本命题是单圈的显式陈述，只涉及 Feynman 积分。设强子 $A$ 在顶点处分裂为活动部分子 $a$（动量 $l$，$l^+>0$）与旁观者 $s_A$（末态动量 $p'$，在壳 $p'^2=m_s^2$，$p'^+>0$），$l=P_A-p'$。一个动量为 $k$ 的胶子附着于 $s_A$，按标准路径 $k$ 经 $s_A$ 进入 $J_A$、经 $a$ 流向硬顶点，于是 $s_A$ 在附着点之前的一段携带 $p'-k$，$a$ 携带 $l+k$。设 $k$ 具有 Glauber 标度，$k^\pm\sim\lambda^2Q$，$|\mathbf k_T|\sim\lambda Q$，并且 $l^-,p'^-,|\mathbf l_T|^2/Q,\dots$ 具有 $J_A$ 的共线标度。则：
\begin{enumerate}[label=(\alph*),leftmargin=2.2em,itemsep=1pt,topsep=2pt]
\item $D_s\equiv(p'-k)^2-m_s^2+\ii0=-2p'^+(k^--c)+\ii0+\order{\lambda^4Q^2}$，$D_a\equiv(l+k)^2-m_a^2+\ii0=2l^+(k^--d)+\ii0+\order{\lambda^4Q^2}$，其中
$c=\frac{2\mathbf p'_T\cdot\mathbf k_T-\mathbf k_T^2}{2p'^+}$，$d=\frac{(\mathbf l_T+\mathbf k_T)^2+m_a^2-2l^+l^-}{2l^+}$ 为实数，均为 $\order{\lambda^2Q}$；$D_s$ 的零点 $k^-=c+\ii0$ 在上半平面，$D_a$ 的零点 $k^-=d-\ii0$ 在下半平面；
\item 在领头幂次（其余因子不依赖 $k^-$）下，$\displaystyle\int\frac{\dd k^-}{2\pi}\,\frac{\ii}{D_s}\,\frac{\ii}{D_a}=\frac{\ii}{4p'^+l^+(c-d)+\ii0}$，并且把积分围道整体移到 $\Imag k^-=\kappa$，只要 $|\kappa|$ 大于两个零点的虚部（$|\kappa|>0$ 的任何固定值即可），积分值都变为 $0$：围道被钳在 $|k^-|\lesssim|c-d|\sim\mathbf k_T^2/p'^+$ 处，不能离开 Glauber 区；
\item 若胶子的另一端附着于 $B$ 的旁观者 $s_B$（旁观者—旁观者交换），则对 $J_B$ 中的 $s_B$ 与活动部分子 $b$ 有同样的结论，其中 $+\leftrightarrow-$、$k\to-k$：$k^+$ 的两个零点分居实轴两侧，$k^+$ 也被钳住；两个分量同时被钳住，即 Glauber 钳夹。
\end{enumerate}
\uses{prop:ir-glauber-trapping,def:ir-glauber}
\end{proposition}

\begin{proof}
(a) 用 $p'^2=m_s^2$ 与 $v\cdot w=v^+w^-+v^-w^+-\mathbf v_T\cdot\mathbf w_T$，
$(p'-k)^2-m_s^2=-2p'^+k^--2p'^-k^++2\mathbf p'_T\cdot\mathbf k_T+2k^+k^--\mathbf k_T^2$。按 Glauber 与共线标度，$2p'^+k^-$、$2\mathbf p'_T\cdot\mathbf k_T$、$\mathbf k_T^2$ 都是 $\order{\lambda^2Q^2}$，而 $2p'^-k^+$、$2k^+k^-$ 是 $\order{\lambda^4Q^2}$；保留前三项即 $-2p'^+(k^--c)$。同样，$(l+k)^2-m_a^2=(2l^+l^--\mathbf l_T^2-m_a^2)+2l^+k^-+2l^-k^+-2\mathbf l_T\cdot\mathbf k_T+2k^+k^--\mathbf k_T^2=2l^+(k^--d)+\order{\lambda^4Q^2}$。零点位置由 $p'^+>0$、$l^+>0$ 与 $+\ii0$ 读出。
(b) 被积函数随 $|k^-|\to\infty$ 衰减为 $|k^-|^{-2}$，可在上半平面闭合围道，只包围 $k^-=c+\ii0$：
\[
\int\frac{\dd k^-}{2\pi}\frac{\ii}{-2p'^+(k^--c)+\ii0}\,\frac{\ii}{2l^+(k^--d)+\ii0}
=\ii\cdot\frac{\ii}{-2p'^+}\cdot\frac{\ii}{2l^+(c-d)+\ii0}=\frac{\ii}{4p'^+l^+(c-d)+\ii0}.
\]
若把实轴整体上移到 $\Imag k^-=\kappa>0$（$\kappa$ 大于零点 $c+\ii0$ 的虚部），新的围道上方没有奇点，在上半平面闭合给出 $0$；若下移到 $\Imag k^-=-\kappa$，在下半平面闭合同样给出 $0$。两种形变都改变了积分值，因此不存在离开实轴附近的形变：被积函数的全部贡献来自两零点之间 $|k^--c|\lesssim|c-d|$ 的区域，而 $|c-d|\sim\mathbf k_T^2/p'^+$。这正是\cref{prop:ir-glauber-trapping}中一条入射线与一条出射线的钳夹判据。
(c) 对 $J_B$，$s_B$ 在发射前携带 $p'_B+k$，$(p'_B+k)^2-m^2=2p'^-_B(k^+-c_B)+\order{\lambda^4Q^2}$，零点 $k^+=c_B-\ii0$；$b$ 携带 $l_b-k$，$(l_b-k)^2-m_b^2=-2l_b^-(k^+-d_B)+\order{\lambda^4Q^2}$，零点 $k^+=d_B+\ii0$。论证同 (b)。
\end{proof}

我们用数值积分核对了 (b)：取 $p'^+=0.6$、$l^+=0.4$、$c=0.013$、$d=-0.021$（任意单位），在实轴上以 $\ii\varepsilon$ 代替 $\ii0$ 求积分，与 $\ii/(4p'^+l^+(c-d))$ 的相对偏差在 $\varepsilon=10^{-4}$、$10^{-6}$ 时分别为 $6\times10^{-3}$、$6\times10^{-5}$，即 $\order{\varepsilon/|c-d|}$，随 $\varepsilon\to0$ 线性趋于零；而把围道移到 $\Imag k^-=\pm0.05$ 时积分在 $10^{-32}$ 的精度内为零。还应注意 (a) 的另一面：在 Glauber 区里，$D_a$ 中被程函近似（只保留 $2l^+k^-$）丢掉的 $-2\mathbf l_T\cdot\mathbf k_T-\mathbf k_T^2$ 与保留的项同阶，程函近似的误差是 $\order{1}$（\cref{def:ir-glauber}）。所以，在实轴上软近似失效，而在实轴附近又无法形变：这两件事合在一起，就是 DY 因子化证明必须专门处理的“Glauber 问题”。
